\cleardoublepage{}
\vspace*{-2em}
\begin{center}
    \bfseries \zihao{3} 致~~~~谢
\end{center}

首先感谢我的导师郑文庭教授的悉心指导。在整个研究过程中，郑老师始终给予我专业的指引，使我受益匪浅。

感谢课题组的两位学长杨凯旋和雷骆，在科研过程中给予我宝贵的建议。

感谢课题组的其他老师和同学们，在交流讨论中，我获得了许多有益的启示。

感谢我的班主任林晓老师在学习与生活上的关心和支持，多次在我迷茫时为我指明方向。

感谢始终尽职尽责、不辞辛苦的辅导员卢佳颖老师，陪伴了我们四年的本科时光。

感谢本科阶段的同学们：方致远、李佳阳、牟阔、王筠凇、郭桂全、陈钧浩等，感谢你们与我一路同行，相互帮助，共同进步，让我度过了充实而难忘的大学时光。

最后，感谢我的家人一直以来的理解与支持，给予我坚实的后盾。

谨以此文，向所有在我成长道路上给予帮助的人致以最诚挚的谢意！

\newpage
\vspace*{-2em}
\begin{center}
    \bfseries \zihao{3} 人工智能工具使用说明
\end{center}

本人郑重声明：本毕业论文为本人独立完成。

在论文撰写过程中，使用了人工智能工具作为辅助，具体情况说明如下：

\begin{itemize}
    \item 使用 GPT-4o（OpenAI，使用日期：2025-05-06）对项目背景和意义部分进行了语句润色与逻辑优化；

    \item 使用 GPT-4o（OpenAI，使用日期：2025-05-06）对同行研发情况部分进行了语句润色与逻辑优化。
\end{itemize}